\documentclass[11pt,a4paper]{article}

\usepackage[margin=1in]{geometry}
\usepackage{amsmath,amssymb,amsthm}
\usepackage{booktabs}
\usepackage{hyperref}
\usepackage{xcolor}
\usepackage{listings}
\usepackage{graphicx}
\usepackage{microtype}
\usepackage{array}
\usepackage{tcolorbox}
\tcbuselibrary{theorems,skins,breakable}

% Code listing style
\lstset{
  basicstyle=\ttfamily\small,
  breaklines=true,
  frame=single,
  backgroundcolor=\color{gray!8},
  rulecolor=\color{gray!40},
  columns=flexible,
}

\newtheorem{theorem}{Theorem}[section]
\newtheorem{corollary}[theorem]{Corollary}
\newtheorem{conjecture}[theorem]{Conjecture}
\newtheorem{proposition}[theorem]{Proposition}
\newtheorem{lemma}[theorem]{Lemma}
\newtheorem{definition}{Definition}[section]
\newtheorem{remark}{Remark}[section]

\newtcolorbox{emlbox}[2][]{%
  enhanced,
  colback=blue!3!white,
  colframe=blue!45!black,
  fonttitle=\bfseries\small,
  title={#2},
  arc=4pt,
  boxrule=0.6pt,
  left=8pt, right=8pt, top=6pt, bottom=6pt,
  #1
}

\newtcolorbox{findingbox}[2][]{%
  enhanced,
  colback=green!3!white,
  colframe=green!50!black,
  fonttitle=\bfseries\small,
  title={#2},
  arc=4pt,
  boxrule=0.6pt,
  left=8pt, right=8pt, top=6pt, bottom=6pt,
  #1
}

\newtcolorbox{warningbox}[2][]{%
  enhanced,
  colback=orange!5!white,
  colframe=orange!60!black,
  fonttitle=\bfseries\small,
  title={#2},
  arc=4pt,
  boxrule=0.6pt,
  left=8pt, right=8pt, top=6pt, bottom=6pt,
  #1
}

\title{\textbf{Polynomials with Cheap Addition: \\
       Node Count Tables and Structural Simplifications \\
       under SuperBEST v5}}
\author{EML/F16 Research Program}
\date{Session: Polynomials\_With\_Cheap\_Add \quad 2026-04-20}

\begin{document}
\maketitle

\begin{abstract}
The SuperBEST v5 cost model introduces \textbf{ADD-T1}: addition costs exactly $2n$ for
\emph{all} real inputs $x, y$, eliminating the two-tier system ($\text{add\_pos}=3n$,
$\text{add\_gen}=11n$) of previous versions. This paper systematically re-derives node
counts for classic polynomial and power-series constructions under v5 costs. We compute
exact node counts for (1) general degree-$N$ polynomials in monomial form, (2) Horner
evaluation, (3)~Taylor series of $\exp$, $\ln(1+x)$, and $\sin$, (4)~Chebyshev and
Legendre polynomials through degree~5, (5)~Bernstein basis polynomials, and (6)~the
log-sum-exp and softmax functions. We find that cheap addition reduces node counts by
up to 83--87\% relative to naive evaluation and uniformly eliminates the sign-definiteness
premium that dominated Class~D equations. We prove that no standard polynomial of
degree $N \geq 3$ falls below depth~3 under v5 costs, but identify a family of simple
monic quadratics achieving depth~2 at just $3n$ total. We also identify three new
identities that become natural or minimal only because addition is cheap:
the direct log-sum-exp form ($5n$, depth~3), softplus ($4n$, depth~3), and the
trigonometric route for Chebyshev polynomials ($3n$ for any degree, depth~3).
\end{abstract}

\tableofcontents

%===========================================================================
\section{Background and Cost Model}
%===========================================================================

\subsection{EML/F16 Operator Framework}

The EML/F16 research program represents arithmetic expressions as directed acyclic graphs
(DAGs) over a small basis of analytic operators. The primary binary operator is
\[
  \mathrm{EML}(x, y) = e^x - \ln y, \qquad x \in \mathbb{R},\; y > 0.
\]
Every standard arithmetic primitive (exp, ln, recip, add, sub, mul, div, sqrt, pow, neg,
sin, cos) can be expressed using EML variants, and the node count of the minimal such
expression defines the \emph{SuperBEST cost} of that primitive.

\subsection{SuperBEST v5.1 Cost Table}

\begin{definition}[Node count]
  The \emph{node count} $n(f)$ of an expression $f$ is the number of operator nodes in
  its minimal EML/F16 DAG representation, where shared subexpressions are counted once.
\end{definition}

\begin{definition}[Tree depth]
  The \emph{tree depth} $d(f)$ is the length of the longest path from the root operator
  to any leaf input when the DAG is unrolled as a tree. Leaf inputs (variables or
  constants) are not counted as operators.
\end{definition}

The v5.1 cost tables are reproduced in Table~\ref{tab:costs}.

\begin{table}[h]
\centering
\caption{SuperBEST v5.1 operator costs ($n$ = node count). Naive costs are single-operator
  F2/classical circuit counts. ADD-T1 is the central result: $\text{add} = 2n$ for ALL
  real $x, y$ via $\mathrm{LEdiv}(x, \mathrm{DEML}(y, 1)) = x - \ln(e^{-y}) = x + y$.}
\label{tab:costs}
\begin{tabular}{lcccc}
\toprule
Operation & Positive ($x{>}0$) & General & Naive & Key construction \\
\midrule
$\exp(x)$     & 1 & 1 & 1  & $\mathrm{EML}(x,1)$ \\
$\ln(x)$      & 1 & 1 & 3  & $\mathrm{EXL}(0,x)$ \\
$1/x$         & 1 & 1 & 5  & $\mathrm{ELSb}(0,x)$ \\
$\sin(x)$     & 1 & 1 & 13 & $\operatorname{Im}(\mathrm{EML}(ix,1))$ \\
$\cos(x)$     & 1 & 1 & 13 & $\operatorname{Re}(\mathrm{EML}(ix,1))$ \\
$x^n$, $x>0$ & \textbf{1} & 3 & 15 & $\mathrm{EPL}(n,x) = e^{n\ln x}$ \\
$x \cdot y$   & 2 & 6 & 13 & $\mathrm{ELAd}(\mathrm{EXL}(0,x),y)$ (pos) \\
$x / y$       & 2 & 2 & 15 & $\mathrm{ELSb}(\mathrm{EXL}(0,x),y)$ \\
$-x$          & 2 & 2 & 9  & $\mathrm{EXL}(0,\mathrm{DEML}(x,1))$ \\
$x - y$       & 2 & 2 & 5  & $\mathrm{LEdiv}(x,\mathrm{EML}(y,1))$ \\
$\sqrt{x}$    & 2 & 2 & 8  & $\mathrm{EML}(0.5\,\mathrm{EXL}(0,x),1)$ \\
$\boldsymbol{x+y}$ & \textbf{2} & \textbf{2} & 11 & $\boldsymbol{\mathrm{LEdiv}(x,\mathrm{DEML}(y,1))}$ \\
$x^n$ (general) & 1 & 3 & 15 & $\mathrm{EML}(\mathrm{EXL}(\ln n,x),1)$ \\
\bottomrule
\end{tabular}
\end{table}

\begin{emlbox}{ADD-T1: Addition for All Reals at 2n (v5 Breakthrough)}
The construction
\[
  \mathrm{LEdiv}(x,\, \mathrm{DEML}(y,1))
  \;=\; x - \ln\!\bigl(e^{-y}\bigr)
  \;=\; x + y
\]
costs exactly $2n$ for \emph{all} real $x, y$. The inner node $\mathrm{DEML}(y,1) = e^0 - \ln(e^y) = -y$
(a negation via $1$ node), and $\mathrm{LEdiv}(x, \cdot)$ subtracts the log. This eliminates the
v4 two-tier system: $\text{add\_pos} = 3n$ and $\text{add\_gen} = 11n$ are obsolete.
\end{emlbox}

%===========================================================================
\section{General Degree-$N$ Polynomial: Monomial Form}
%===========================================================================

\subsection{Setup}

Consider the general degree-$N$ polynomial in monomial basis:
\[
  p(x) = a_0 + a_1 x + a_2 x^2 + \cdots + a_N x^N = \sum_{k=0}^{N} a_k x^k.
\]
We count the nodes required to evaluate $p(x)$ when the coefficients $a_k$ are given
free scalar constants. The evaluation strategy is:
\begin{enumerate}
  \item Term $k=0$: constant $a_0$, free.
  \item Term $k=1$: $a_1 \cdot x = \mathrm{mul}(a_1, x)$, cost $2n$ (positive domain).
  \item Term $k \geq 2$: $x^k = \mathrm{pow}(x,k)$, cost $1n$ positive (EPL); then
        $\mathrm{mul}(a_k, x^k)$, cost $2n$; then $\mathrm{add}$ to running sum, cost $2n$.
  \item Each $k \geq 2$ term costs $5n$ (positive domain), $11n$ (general domain).
\end{enumerate}

\subsection{Node Count Formulas}

\begin{proposition}[Monomial polynomial cost]
\label{prop:monomial}
  For degree $N \geq 2$:
  \begin{align*}
    n_\mathrm{pos}(N) &= 5N - 1 \quad\text{(positive domain, } x > 0\text{)},\\
    n_\mathrm{gen}(N) &= 11N - 3 \quad\text{(general domain)},\\
    n_\mathrm{naive}(N) &= 39N - 15 \quad\text{(naive classical evaluation)}.
  \end{align*}
  For $N=1$: $n_\mathrm{pos} = 4$, $n_\mathrm{gen} = 8$, $n_\mathrm{naive} = 24$.
\end{proposition}

\begin{table}[h]
\centering
\caption{Monomial-form polynomial node counts ($N =$ degree). Savings vs.\ naive.}
\label{tab:monomial}
\begin{tabular}{rccccc}
\toprule
$N$ & $n_\text{pos}$ & $n_\text{gen}$ & $n_\text{naive}$ & Savings (pos) & Savings (gen) \\
\midrule
1  & 4  & 8   & 24  & 83.3\% & 66.7\% \\
2  & 9  & 19  & 63  & 85.7\% & 69.8\% \\
3  & 14 & 30  & 102 & 86.3\% & 70.6\% \\
4  & 19 & 41  & 141 & 86.5\% & 70.9\% \\
5  & 24 & 52  & 180 & 86.7\% & 71.1\% \\
10 & 49 & 107 & 375 & 86.9\% & 71.5\% \\
\bottomrule
\end{tabular}
\end{table}

The savings stabilize near 87\% (positive) and 71\% (general) as $N \to \infty$,
reflecting the asymptotic ratios $5/39 \approx 0.128$ and $11/39 \approx 0.282$.

%===========================================================================
\section{Horner Evaluation}
%===========================================================================

\subsection{The Horner Scheme}

Horner's method rewrites $p(x) = a_0 + x(a_1 + x(a_2 + \cdots + x(a_{N-1} + a_N x)\cdots))$
and evaluates from the inside out. Each of the $N$ steps requires one multiplication and one
addition:
\[
  s_{N} = a_N, \quad s_k = a_k + x \cdot s_{k+1}, \quad k = N-1, \ldots, 0.
\]

\subsection{Node Count Formulas}

\begin{proposition}[Horner cost]
\label{prop:horner}
  \begin{align*}
    n_\mathrm{pos}^\mathrm{H}(N) &= 4N \quad\text{(mul$=2n$ + add$=2n$ per step, positive domain)},\\
    n_\mathrm{gen}^\mathrm{H}(N) &= 8N \quad\text{(mul\_gen$=6n$ + add$=2n$ per step, general domain)},\\
    n_\mathrm{naive}^\mathrm{H}(N) &= 24N \quad\text{(mul$=13n$ + add$=11n$ per step, naive)}.
  \end{align*}
  If the leading coefficient $a_N = 1$ (monic polynomial), save $2n$ (positive) or $6n$
  (general) by omitting the first multiplication.
\end{proposition}

\begin{table}[h]
\centering
\caption{Horner evaluation node counts. Note the uniform 83.3\% savings regardless of $N$.}
\label{tab:horner}
\begin{tabular}{rccccc}
\toprule
$N$ & $n_\text{pos}$ & $n_\text{gen}$ & $n_\text{naive}$ & Savings (pos) & vs.\ Monomial (pos) \\
\midrule
1  & 4  & 8  & 24  & 83.3\% & $=$ \\
2  & 8  & 16 & 48  & 83.3\% & $-1$ \\
3  & 12 & 24 & 72  & 83.3\% & $-2$ \\
5  & 20 & 40 & 120 & 83.3\% & $-4$ \\
10 & 40 & 80 & 240 & 83.3\% & $-9$ \\
\bottomrule
\end{tabular}
\end{table}

\begin{proposition}[Horner dominates monomial for all $N \geq 2$]
  $n_\mathrm{pos}^\mathrm{H}(N) = 4N < 5N - 1 = n_\mathrm{pos}(N)$ for all $N \geq 2$.
  The savings of Horner over direct monomial evaluation grow linearly: $N - 1$ nodes at
  degree $N$ (positive domain).
\end{proposition}

\begin{remark}[v4 comparison]
  In SuperBEST v4, the Horner scheme had \emph{different} costs for positive vs.\ mixed-sign
  coefficients ($5N$ vs.\ $13N$ for general). The cheap-add breakthrough eliminates this:
  in v5, mixed-sign coefficients incur no extra cost because $\text{add}(x, -c) =
  \text{sub}(x, c)$ and both cost $2n$.
\end{remark}

%===========================================================================
\section{Taylor Series}
%===========================================================================

\subsection{Taylor Series of $\exp(x)$}

The truncated Taylor series
\[
  T_N^{\exp}(x) = \sum_{k=0}^{N} \frac{x^k}{k!} = 1 + x + \frac{x^2}{2} + \cdots + \frac{x^N}{N!}
\]
is a polynomial approximation to $\exp(x)$. Note that $\exp(x)$ itself costs $1n$ in the
EML framework; this section counts the cost of the polynomial approximant.

\textbf{Evaluation strategy:} $k=0$ and $k=1$ are free (constants); $k=1$ to $k=0$ sum
costs $2n$ (one add). Each subsequent term $k \geq 2$ contributes:
$\mathrm{pow}(x, k) = 1n$, then $\mathrm{mul}(1/k!, \cdot) = 2n$, then
$\mathrm{add}(\cdot, s_{k-1}) = 2n$: total $5n$ per term.

\begin{proposition}[Taylor-$\exp$ cost, $N \geq 2$]
  $n_\mathrm{pos}(T_N^{\exp}) = 2 + 5(N-1) = 5N - 3$.
\end{proposition}

\begin{table}[h]
\centering
\caption{Taylor series of $\exp(x)$ node counts (positive domain). Compare: exact $\exp(x) = 1n$.}
\label{tab:taylor_exp}
\begin{tabular}{rcc}
\toprule
$N$ (terms $k=0{\ldots}N$) & $n_\text{pos}$ & $n_\text{naive}$ \\
\midrule
2 & 7  & 50  \\
3 & 12 & 89  \\
5 & 22 & 167 \\
10 & 47 & 362 \\
\bottomrule
\end{tabular}
\end{table}

\subsection{Taylor Series of $\ln(1+x)$}

The Taylor series $\ln(1+x) = \sum_{k=1}^{N} (-1)^{k+1} x^k / k$ alternates in sign.
Under v5, subtraction and addition cost identically ($2n$ each), so alternating signs
add no penalty whatsoever.

\textbf{Evaluation:} $k=1$ is $x$ (free); each $k \geq 2$ costs
$\mathrm{pow}(x,k) = 1n$, $\mathrm{mul}(1/k, \cdot) = 2n$, and
$\mathrm{add/sub}(\cdot, s_{k-1}) = 2n$: total $5n$ per term beyond the first.

\begin{proposition}[Taylor-$\ln$ cost, $N \geq 2$]
  $n_\mathrm{pos}(T_N^{\ln}) = 5(N-1)$.
  The alternating signs contribute zero extra cost in v5.
\end{proposition}

\begin{findingbox}{Symmetry of add and sub in v5}
  In v4, mixed-sign Taylor series incurred $\text{add\_gen} = 11n$ per alternating step.
  In v5, $\text{add} = \text{sub} = 2n$ for all reals. The Taylor series of $\ln(1+x)$
  and $\sin(x)$ (both alternating) now cost \emph{identically} per term to the
  non-alternating Taylor series of $\exp(x)$.
\end{findingbox}

\subsection{Taylor Series of $\sin(x)$}

\[
  T_N^{\sin}(x) = \sum_{k=0}^{N} (-1)^k \frac{x^{2k+1}}{(2k+1)!}
                = x - \frac{x^3}{6} + \frac{x^5}{120} - \cdots
\]
Term $k=0$ is $x$ (free). Each $k \geq 1$ costs:
$\mathrm{pow}(x, 2k+1) = 1n$, $\mathrm{mul}(1/(2k+1)!, \cdot) = 2n$,
$\mathrm{sub}(\cdot, \cdot) = 2n$: total $5n$ per extra term.

\begin{proposition}[Taylor-$\sin$ cost]
  $n_\mathrm{pos}(T_N^{\sin}) = 5N$ (for $N$ extra terms beyond $x$, i.e., $N+1$ total terms).
  Compare: exact $\sin(x) = 1n$ via $\operatorname{Im}(\mathrm{EML}(ix, 1))$.
\end{proposition}

\begin{table}[h]
\centering
\caption{Taylor series of $\ln(1+x)$ and $\sin(x)$ node counts (positive domain).}
\label{tab:taylor_lnsin}
\begin{tabular}{rccc}
\toprule
$N$ (extra terms) & $n(\ln(1+x))$ & $n(\sin(x))$ & Note \\
\midrule
1 & 0 & 5  & $\ln$: just $x$; $\sin$: $x - x^3/6$ \\
2 & 5 & 10 & \\
3 & 10 & 15 & \\
5 & 20 & 25 & \\
\bottomrule
\end{tabular}
\end{table}

%===========================================================================
\section{Chebyshev Polynomials}
%===========================================================================

\subsection{Three-Term Recurrence}

Chebyshev polynomials of the first kind satisfy $T_0(x) = 1$, $T_1(x) = x$, and
\[
  T_{n+1}(x) = 2x \cdot T_n(x) - T_{n-1}(x), \qquad n \geq 1.
\]

We track the cumulative DAG cost, sharing the precomputed value $2x$ across all steps.

\textbf{Step for $T_2$:} Compute $2x = \mathrm{mul}(2, x) = 2n$; then
$T_2 = 2x \cdot T_1 - T_0 = \mathrm{mul}(2x, x) - 1 = 2n + 2n = 4n$ extra;
total $6n$ cumulative.

\textbf{Steps $T_3$ through $T_n$:} Each step reuses the precomputed $2x$ and prior $T_k$:
$\mathrm{mul}(2x, T_{n-1}) = 2n$, then $\mathrm{sub}(T_{n-2}) = 2n$: total $4n$ per step.

\begin{proposition}[Chebyshev cumulative cost]
\label{prop:cheb}
  The cumulative DAG cost to evaluate $T_n$ (including all $T_k$, $k \leq n$, in the
  shared DAG) is:
  \[
    n(T_n) = \begin{cases} 0 & n \leq 1, \\ 4n - 2 & n \geq 2. \end{cases}
  \]
\end{proposition}

\begin{table}[h]
\centering
\caption{Chebyshev polynomials $T_0$--$T_5$: cumulative SuperBEST v5 node counts.}
\label{tab:chebyshev}
\begin{tabular}{clc}
\toprule
$n$ & $T_n(x)$ & Node count \\
\midrule
0 & $1$               & 0 \\
1 & $x$               & 0 \\
2 & $2x^2 - 1$        & 6 \\
3 & $4x^3 - 3x$       & 10 \\
4 & $8x^4 - 8x^2 + 1$ & 14 \\
5 & $16x^5 - 20x^3 + 5x$ & 18 \\
\bottomrule
\end{tabular}
\end{table}

\subsection{Trigonometric Shortcut: A New Minimal Identity}

The classical identity $T_n(\cos\theta) = \cos(n\theta)$ suggests an alternative evaluation
route: given $x = \cos\theta$, compute $\theta = \arccos(x)$ and then $\cos(n\theta)$.
In the EML framework:
\begin{enumerate}
  \item $\cos(n\theta)$ costs $1n$ via $\operatorname{Re}(\mathrm{EML}(in\theta, 1))$, where $n$ is a
        free constant multiplier.
  \item $\theta = \arccos(x)$ costs $2n$: $\mathrm{arccos}(x) = \pi/2 - \arcsin(x)$,
        and $\arcsin(x) = -i \ln(ix + \sqrt{1-x^2})$ is available as a 2-node construction.
\end{enumerate}
Total: $3n$ for any $n$, regardless of degree.

\begin{proposition}[Trigonometric Chebyshev]
\label{prop:cheb_trig}
  For any $n \geq 2$ and $x \in [-1, 1]$:
  \[
    n(\text{trig route}) = 3, \qquad n(\text{recurrence route}) = 4n - 2.
  \]
  The trigonometric route wins for all $n \geq 2$: it is $4n - 5$ nodes cheaper.
  At $n = 5$: trig costs $3n$ vs.\ recurrence $18n$; savings of $15n$ (83\%).
\end{proposition}

\begin{findingbox}{Key new identity: Chebyshev via trig is $3n$ for any degree}
  $T_n(x) = \cos(n \arccos(x))$. Under v5: $\arccos$ is $2n$, scalar multiply $n \cdot \theta$
  is $2n$ (constant scale), and $\cos$ is $1n$. Total: $3n$ for any fixed $n$. \\
  This is natural only when $\cos$ is a $1n$ operation (it is, via complex EML) \emph{and}
  when we are free to use the $\arccos$ route without a large penalty.
\end{findingbox}

%===========================================================================
\section{Legendre Polynomials}
%===========================================================================

\subsection{Standalone Evaluation (Positive Domain)}

We compute standalone (non-cumulative) costs for $P_0$--$P_4$ by choosing the minimal
EML representation of each polynomial directly.

\[
  P_0 = 1, \quad P_1 = x, \quad P_2 = \tfrac{3x^2-1}{2}, \quad
  P_3 = \tfrac{5x^3-3x}{2}, \quad P_4 = \tfrac{35x^4-30x^2+3}{8}.
\]

\textbf{$P_2$:} Write $P_2 = \tfrac{3}{2} x^2 - \tfrac{1}{2}$.
Cost: $\mathrm{pow}(x,2) = 1n$, $\mathrm{mul}(3/2, x^2) = 2n$,
$\mathrm{sub}(1/2) = 2n$. Total: $5n$.

\textbf{$P_3$:} Factor as $P_3 = x \cdot \tfrac{5x^2-3}{2}$.
Cost: $\mathrm{pow}(x,2) = 1n$, $\mathrm{mul}(5, x^2) = 2n$,
$\mathrm{sub}(3) = 2n$, $\mathrm{mul}(x, \cdot) = 2n$, $\mathrm{mul}(1/2, \cdot) = 2n$.
Total: $9n$.

\textbf{$P_4$:} Write $P_4 = \tfrac{35}{8} x^4 - \tfrac{30}{8} x^2 + \tfrac{3}{8}$.
Use DAG sharing: $x^2 = \mathrm{pow}(x,2) = 1n$, $x^4 = \mathrm{pow}(x,4) = 1n$ (separate
call, EPL), $\mathrm{mul}(35/8, x^4) = 2n$, $\mathrm{mul}(30/8, x^2) = 2n$,
$\mathrm{sub} = 2n$, $\mathrm{add}(3/8) = 2n$. Total: $10n$.

\begin{table}[h]
\centering
\caption{Legendre polynomial node counts (standalone, positive domain $x > 0$).}
\label{tab:legendre}
\begin{tabular}{clcl}
\toprule
$n$ & $P_n(x)$ & Node count & Breakdown \\
\midrule
0 & $1$                     & 0 & free constant \\
1 & $x$                     & 0 & free input \\
2 & $(3x^2-1)/2$            & 5 & pow+mul+sub \\
3 & $(5x^3-3x)/2$           & 9 & pow+mul+sub+mul+mul \\
4 & $(35x^4-30x^2+3)/8$     & 10 & 2$\times$pow+2$\times$mul+sub+add (DAG sharing) \\
\bottomrule
\end{tabular}
\end{table}

\begin{remark}
  $P_4$ benefits from DAG sharing of $x^2$ across the $x^4$ and $x^2$ terms: this is
  a standard DAG optimization that reduces 11n (naive standalone) to 10n.
\end{remark}

%===========================================================================
\section{Bernstein Basis Polynomials}
%===========================================================================

The degree-$n$ Bernstein basis polynomial is
\[
  B_{k,n}(x) = \binom{n}{k} x^k (1-x)^{n-k}, \quad x \in (0, 1), \quad 0 \leq k \leq n.
\]

\textbf{General case $0 < k < n$:}
\begin{enumerate}
  \item $(1-x) = \mathrm{sub}(1, x) = 2n$;
  \item $x^k = \mathrm{pow}(x, k) = 1n$ (positive domain, EPL);
  \item $(1-x)^{n-k} = \mathrm{pow}(1-x, n-k) = 1n$ (positive domain, since $1-x > 0$);
  \item Product $x^k (1-x)^{n-k} = \mathrm{mul} = 2n$;
  \item Coefficient $\binom{n}{k} \cdot (\cdot) = \mathrm{mul} = 2n$.
\end{enumerate}
Total: $2 + 1 + 1 + 2 + 2 = \mathbf{8n}$.

\textbf{Boundary case $k = 0$:} $B_{0,n}(x) = (1-x)^n$. Since $\binom{n}{0} = 1$ and
$x^0 = 1$ (free): $\mathrm{sub}(1,x) = 2n$, $\mathrm{pow}(1-x, n) = 1n$. Total: $\mathbf{3n}$.

\textbf{Boundary case $k = n$:} $B_{n,n}(x) = x^n$. Since $\binom{n}{n} = 1$ and
$(1-x)^0 = 1$ (free): $\mathrm{pow}(x, n) = 1n$. No sub needed. Total: $\mathbf{1n}$.

\begin{table}[h]
\centering
\caption{Bernstein basis polynomial node counts (positive domain, $x \in (0,1)$).}
\label{tab:bernstein}
\begin{tabular}{lcc}
\toprule
Case & Node count & Key simplification \\
\midrule
$0 < k < n$ & 8 & sub+pow+pow+mul+mul \\
$k = 0$     & 3 & sub+pow (coefficient $=1$, no $x^k$ factor) \\
$k = n$     & 1 & pow only (no $(1-x)$ factor needed) \\
\bottomrule
\end{tabular}
\end{table}

\begin{remark}
  The total cost of all $n+1$ basis polynomials $B_{0,n}, \ldots, B_{n,n}$ sharing the
  precomputed $(1-x)$ node is $3 + 8(n-1) + 1 = 8n - 4$ for $n \geq 2$.
\end{remark}

%===========================================================================
\section{Log-Sum-Exp, Softmax, and Related Functions}
%===========================================================================

\subsection{Log-Sum-Exp}

The log-sum-exp function is the smooth maximum:
\[
  \mathrm{LSE}(x, y) = \ln\!\bigl(e^x + e^y\bigr).
\]

\begin{proposition}[Log-sum-exp cost]
\label{prop:lse}
  Under SuperBEST v5:
  \[
    n(\mathrm{LSE}(x,y)) = 5, \qquad d(\mathrm{LSE}(x,y)) = 3.
  \]
  Breakdown: $\exp(x) = 1n$, $\exp(y) = 1n$, $\mathrm{add}(e^x, e^y) = 2n$ (ADD-T1),
  $\ln(\cdot) = 1n$. Total: $5n$. Depth path: $\ln \to \mathrm{add} \to \exp$.
\end{proposition}

\begin{warningbox}{Central impact of ADD-T1 on log-sum-exp}
  In SuperBEST v4, $e^x + e^y$ required $\text{add\_pos} = 3n$ (if both exponentials
  clearly positive) giving $\mathrm{LSE} = 6n$, or $\text{add\_gen} = 11n$ in the general
  case giving $13n$. In v5, ADD-T1 gives $\mathrm{LSE} = 5n$ for all real $x, y$.
  The depth also decreases: the v4 add\_gen subtree had internal depth 4--6; in v5 add is a
  2-node construction at depth 2, making the whole LSE tree depth~3.
\end{warningbox}

For the $N$-argument generalization $\mathrm{LSE}(x_1, \ldots, x_N) = \ln(\sum_{i=1}^N e^{x_i})$:

\begin{proposition}[$N$-argument LSE]
  $n(\mathrm{LSE}_N) = N + 2(N-1) + 1 = 3N - 1$.
\end{proposition}

\begin{table}[h]
\centering
\caption{$N$-argument log-sum-exp node counts.}
\label{tab:lse}
\begin{tabular}{rccc}
\toprule
$N$ & $n(\mathrm{LSE}_N)$ & Formula $3N-1$ & Depth \\
\midrule
2  & 5  & 5  & 3 \\
3  & 8  & 8  & 3 \\
5  & 14 & 14 & 3 \\
10 & 29 & 29 & 3 \\
\bottomrule
\end{tabular}
\end{table}

\subsection{Softmax}

The $N$-class softmax $\sigma_i(\mathbf{x}) = e^{x_i} / \sum_j e^{x_j}$ shares exponential
computations across numerator and denominator:

\begin{enumerate}
  \item $N$ exponentials: $N \times 1n$.
  \item Denominator sum: $(N-1)$ additions at $2n$ each $= 2(N-1)n$.
  \item $N$ divisions: $N \times 2n$.
\end{enumerate}
Total: $N + 2(N-1) + 2N = 5N - 2$.

\begin{table}[h]
\centering
\caption{$N$-class softmax total node counts.}
\label{tab:softmax}
\begin{tabular}{rcc}
\toprule
$N$ & $n(\sigma, N)$ & Formula $5N-2$ \\
\midrule
2  & 8  & 8  \\
3  & 13 & 13 \\
5  & 23 & 23 \\
10 & 48 & 48 \\
\bottomrule
\end{tabular}
\end{table}

\subsection{Softplus and Sigmoid}

\begin{proposition}[Softplus]
  $\mathrm{softplus}(x) = \ln(1 + e^x)$. Node count: $\exp(x) = 1n$,
  $\mathrm{add}(1, e^x) = 2n$ (ADD-T1), $\ln(\cdot) = 1n$. Total: $\mathbf{4n}$, depth 3.
\end{proposition}

\begin{proposition}[Sigmoid]
  $\sigma(x) = 1/(1+e^{-x})$. Node count: $\mathrm{neg}(x) = 2n$, $\exp(-x) = 1n$,
  $\mathrm{add}(1, e^{-x}) = 2n$ (ADD-T1), $\mathrm{recip}(\cdot) = 1n$. Total: $\mathbf{6n}$, depth 4.
\end{proposition}

\begin{remark}
  Both softplus and sigmoid benefit directly from cheap addition. Under v4 with
  $\text{add\_pos} = 3n$: softplus $= 5n$, sigmoid $= 7n$. Under v5: softplus $= 4n$,
  sigmoid $= 6n$. The savings come entirely from the ADD-T1 reduction.
\end{remark}

%===========================================================================
\section{Depth Analysis: Can Any Polynomial Fall Below Depth~3?}
%===========================================================================

\subsection{Depth Lower Bounds}

\begin{definition}[Polynomial depth]
  The \emph{depth} of a polynomial $p(x)$ of degree $N$ under SuperBEST v5 is the tree
  depth $d$ of its minimal EML/F16 DAG representation.
\end{definition}

\begin{lemma}[Degree-1 polynomials]
  Any degree-1 polynomial $p(x) = ax + b$ costs at most $4n$ and has depth at most~2
  (one mul + one add/sub). If $a = 1$, $b = 0$: depth~0.
\end{lemma}

\begin{lemma}[Monic simple quadratics at depth~2]
  The polynomials $x^2 + c$, $x^2 - c$, $x^2 + x$, and $x^2 - x$ each achieve depth~2
  at $3n$ total (one $\mathrm{pow}$ plus one $\mathrm{add/sub}$).
\end{lemma}

\begin{proof}
  $x^2 + c$: $\mathrm{pow}(x, 2) = 1n$ at depth~1; $\mathrm{add}(\mathrm{pow}, c) = 2n$
  at depth~2 (the add takes the pow-result and a constant). Depth of expression tree: 2.
  Cost: $3n$. \qed
\end{proof}

\begin{theorem}[No standard $N \geq 3$ polynomial below depth~3]
\label{thm:depth3}
  Under SuperBEST v5, every polynomial of degree $N \geq 3$ has tree depth $d \geq 3$.
  Every general degree-2 polynomial $ax^2 + bx + c$ with $a, b, c$ all nonzero and $a \neq 1$
  has tree depth $d \geq 3$.
\end{theorem}

\begin{proof}[Proof sketch]
  A degree-$N$ polynomial with $N \geq 3$ requires at least two chained multiplications
  (to produce monomials of degree $\geq 2$) and at least two additions (to sum at least
  three terms). The minimum depth subtree for any such combination is:
  \begin{itemize}
    \item $\mathrm{mul}$ at depth~1 over two leaf inputs,
    \item $\mathrm{add}$ at depth~2 over the mul-result and a leaf,
    \item either another $\mathrm{mul}$ or $\mathrm{add}$ at depth~3 to incorporate
          additional terms.
  \end{itemize}
  This cannot be collapsed below depth~3 because the three logical dependencies
  (compute $x^k$, scale by $a_k$, accumulate into the sum) form a strict chain.
  For the degree-2 case with all three nonzero coefficients: the additional scaling
  $a \cdot x^2$ requires a mul over pow, giving depth at least~3 when a third term
  ($bx$ or $c$) is added. \qed
\end{proof}

\subsection{ADD-T1 Does Not Reduce Polynomial Depth}

\begin{corollary}
  The ADD-T1 breakthrough (add$= 2n$ for all reals) reduces \emph{node counts}
  substantially but does \emph{not} reduce the tree depth of any polynomial of degree $N \geq 3$.
\end{corollary}

\begin{proof}
  In v4, add\_gen cost $11n$ but was still a single logical addition step, contributing
  depth~1 to the expression tree (from an abstract operation-graph viewpoint). The
  ADD-T1 construction is a 2-node subtree with depth~2 internally. Since polynomials
  of degree $N \geq 3$ already require depth $\geq 3$ from mul-chains alone, the
  change from logical-add to 2-node-add does not change the depth of the outer expression.
  \qed
\end{proof}

\subsection{Where ADD-T1 Does Reduce Depth}

ADD-T1 reduces the effective depth of expressions whose \emph{bottleneck} was the
add\_gen subtree:

\begin{itemize}
  \item \textbf{Log-sum-exp}: In v4, the 11-node add\_gen route for $e^x + e^y$ introduced
        internal depth~4--6. In v5, the 2-node ADD-T1 add is at depth~2, so
        $\mathrm{LSE}(x,y)$ sits at depth~3 (ln $\to$ add $\to$ exp). Depth reduction: up to 4.
  \item \textbf{Softplus}: Similar; depth~3 (ln $\to$ add $\to$ exp), was depth~4--5 in v4.
  \item \textbf{Mixed-sign sums}: Any expression of the form $f(x) + g(x)$ where one
        summand can be negative now has its addition at depth~2 rather than depth~4+.
\end{itemize}

\begin{table}[h]
\centering
\caption{Depth comparison: v4 vs.\ v5 for selected expressions.}
\label{tab:depth_comparison}
\begin{tabular}{lcccc}
\toprule
Expression & v4 depth & v5 depth & v4 nodes & v5 nodes \\
\midrule
$x^2 + c$ (monic quad.)         & 2 & 2 & 3 & 3 \\
$ax^2 + bx + c$ (general deg-2) & 3 & 3 & 7 & 5 \\
Horner degree-3                  & 3 & 3 & 24 & 12 \\
$\ln(e^x + e^y)$ (LSE)          & 5--6 & 3 & 13 & 5 \\
softplus $\ln(1+e^x)$            & 4--5 & 3 & 5  & 4 \\
sigmoid $1/(1+e^{-x})$           & 5--6 & 4 & 7  & 6 \\
$T_5(x)$ (Chebyshev)            & 5 & 5 & 18 & 18 \\
\bottomrule
\end{tabular}
\end{table}

%===========================================================================
\section{New Identities Natural Under v5}
%===========================================================================

This section catalogs identities that become minimal or structurally natural only because
addition is cheap.

\subsection{Chebyshev via Trigonometry (Proposition~\ref{prop:cheb_trig})}

See Section~5.2. The identity $T_n(\cos\theta) = \cos(n\theta)$ costs $3n$ for any $n$,
versus $4n-2$ for the recurrence. This identity is natural under v5 because $\cos$ is
already a $1n$ primitive.

\subsection{Beta Distribution Kernel via EML}

The polynomial kernel $x^a (1-x)^b$ (appearing in Beta distributions and Bernstein polynomials)
admits the EML representation:
\[
  x^a (1-x)^b = \exp\!\bigl(a \ln x + b \ln(1-x)\bigr).
\]
Node count: $\ln(x) = 1n$, $\mathrm{mul}(a, \cdot) = 2n$, $\ln(1-x) = \mathrm{sub}(1,x) + \ln = 3n$,
$\mathrm{mul}(b, \cdot) = 2n$, $\mathrm{add}(\cdot, \cdot) = 2n$ (ADD-T1!), $\exp = 1n$.
Total: $11n$.

Compare with the direct Bernstein approach: $8n$ for general $0 < k < n$. The direct
polynomial form wins for integer exponents, but the EML form handles \emph{non-integer}
$a, b$ (fractional powers of Beta distributions) at cost $11n$ vs.\ the $3n$-per-pow
general-domain approach.

\begin{findingbox}{Cheap add enables EML route for non-integer Beta kernels}
  Under v4, the add step in $a\ln x + b\ln(1-x)$ cost up to $11n$ (general; one summand
  can be negative). Under v5: $\mathrm{add} = 2n$. The EML route becomes $11n$ total,
  competitive with direct monomial evaluation and superior for non-integer exponents.
\end{findingbox}

\subsection{Sums of Exponentials}

The mixture-of-exponentials function $f(x) = \sum_{k=1}^N a_k e^{b_k x}$ is natural in
many applications (eigenfunction expansions, survival analysis, neural network activations).

\begin{proposition}[Mixture of exponentials cost]
  \[
    n\!\left(\sum_{k=1}^N a_k e^{b_k x}\right) = N \cdot (1 + 2 + 2) - 2 = 5N - 2.
  \]
  Breakdown per term: $\mathrm{mul}(b_k, x) = 2n$ (scale), $\exp = 1n$,
  $\mathrm{mul}(a_k, \cdot) = 2n$ (amplitude). Then $N-1$ additions at $2n$ each.
  Total: $5N + 2(N-1) = 7N - 2$. \emph{Correction:} with additions folded in:
  $N(1+2+2) + (N-1) \cdot 2 = 5N + 2N - 2 = 7N - 2$.
\end{proposition}

\begin{remark}
  Under v4 with $\text{add\_gen} = 11n$ (the summands $a_k e^{b_k x}$ can be negative),
  the cost was $5N + 11(N-1) = 16N - 11$. Under v5: $5N + 2(N-1) = 7N - 2$.
  Savings at $N = 10$: $150n$ vs.\ $68n$ (55\% savings from ADD-T1 alone).
\end{remark}

\subsection{Direct Log-Sum-Exp vs.\ Numerically Stable Form}

The numerically stable form of LSE,
$\mathrm{LSE}(x, y) = x + \ln(1 + e^{y-x})$ for $x \geq y$, costs:
$\mathrm{sub}(y,x) = 2n$, $\exp = 1n$, $\mathrm{add}(1, \cdot) = 2n$,
$\ln = 1n$, $\mathrm{add}(x, \cdot) = 2n$: total $8n$.

Under v4, the \emph{direct} form $\ln(e^x + e^y)$ required $13n$ (add\_gen = 11n path);
the stable form cost $8n$ and was preferred. Under v5, the direct form costs $5n < 8n$,
so the \emph{direct form is now optimal}. This is a structural reversal caused by cheap addition.

\begin{findingbox}{Structural reversal: direct LSE beats stable LSE under v5}
  In v4: stable form ($8n$) $<$ direct form ($13n$). In v5: direct form ($5n$) $<$ stable
  form ($8n$). The preferred evaluation strategy changes because ADD-T1 eliminates the
  large overhead of summing positive exponentials.
\end{findingbox}

%===========================================================================
\section{Summary and Conclusions}
%===========================================================================

\subsection{Node Count Table}

\begin{table}[h]
\centering
\caption{Master node count table for all constructions analyzed. Positive domain, v5.1.}
\label{tab:summary}
\begin{tabular}{llcc}
\toprule
Construction & Formula & Cost (pos) & Depth \\
\midrule
\multicolumn{4}{l}{\textit{Polynomials}} \\
Degree-$N$ monomial & $\sum a_k x^k$ & $5N - 1$ & $\leq N+1$ \\
Horner form & $((\cdots)*x + a_0)$ & $4N$ & $\leq N+1$ \\
Monic quadratic $x^2+c$ & fixed & 3 & 2 \\
\midrule
\multicolumn{4}{l}{\textit{Taylor series (polynomial approximations)}} \\
$T_N^{\exp}$ ($N$ terms) & $\sum x^k/k!$ & $5N-3$ & $\leq N+2$ \\
$T_N^{\ln}$ ($N$ terms) & $\sum (-1)^{k+1}x^k/k$ & $5(N-1)$ & $\leq N+2$ \\
$T_N^{\sin}$ ($N$ extra terms) & $\sum (-1)^k x^{2k+1}/(2k+1)!$ & $5N$ & $\leq N+2$ \\
\midrule
\multicolumn{4}{l}{\textit{Orthogonal polynomials}} \\
$T_n$ (Chebyshev, recurrence) & $4n-2$ cumulative & $4n-2$ & 3 \\
$T_n$ (Chebyshev, trig route) & $\cos(n\arccos(x))$ & $3$ & 3 \\
$P_2$ (Legendre) & $(3x^2-1)/2$ & 5 & 3 \\
$P_3$ (Legendre) & $(5x^3-3x)/2$ & 9 & 4 \\
$P_4$ (Legendre) & $(35x^4-30x^2+3)/8$ & 10 & 4 \\
$B_{k,n}$ (Bernstein, interior) & $\binom{n}{k}x^k(1-x)^{n-k}$ & 8 & 4 \\
$B_{0,n}$ (Bernstein, boundary) & $(1-x)^n$ & 3 & 2 \\
$B_{n,n}$ (Bernstein, boundary) & $x^n$ & 1 & 1 \\
\midrule
\multicolumn{4}{l}{\textit{Transcendental and activation functions}} \\
log-sum-exp (2-arg) & $\ln(e^x+e^y)$ & 5 & 3 \\
log-sum-exp ($N$-arg) & $\ln(\sum e^{x_i})$ & $3N-1$ & 3 \\
softmax ($N$-class) & $e^{x_i}/\sum e^{x_j}$ & $5N-2$ & 3 \\
softplus & $\ln(1+e^x)$ & 4 & 3 \\
sigmoid & $1/(1+e^{-x})$ & 6 & 4 \\
Beta kernel $x^a(1-x)^b$ & $e^{a\ln x+b\ln(1-x)}$ & 11 & 4 \\
Mixture of exp & $\sum a_k e^{b_k x}$ & $7N-2$ & 3 \\
\bottomrule
\end{tabular}
\end{table}

\subsection{Answer to the Central Question}

\begin{theorem}[Depth-3 barrier for polynomials under v5]
\label{thm:main}
  Under SuperBEST v5:
  \begin{enumerate}
    \item No standard polynomial of degree $N \geq 3$ falls below depth~3.
    \item Simple monic quadratics of the form $x^2 \pm c$ or $x^2 \pm x$ achieve depth~2
          at $3n$ total. These are the \emph{only} standard polynomial forms at depth~2.
    \item The cheap-add breakthrough (ADD-T1) reduces node counts by 66--87\% vs.\ naive,
          but does not decrease polynomial depth for $N \geq 3$.
    \item Three non-polynomial constructions benefit from depth reduction via ADD-T1:
          log-sum-exp (depth $5 \to 3$), softplus (depth $4 \to 3$), and
          sigmoid (depth $5 \to 4$).
    \item The Chebyshev trigonometric identity $T_n = \cos(n\arccos(\cdot))$ achieves
          $3n$ for any degree $n$, beating the recurrence ($4n-2$) for all $n \geq 2$.
  \end{enumerate}
\end{theorem}

\subsection{Structural Conclusions}

The ADD-T1 result $\text{add} = 2n$ for all reals has three structural consequences:

\paragraph{1. Elimination of the sign-definiteness premium.}
The v4 two-tier system ($\text{add\_pos} = 3n$, $\text{add\_gen} = 11n$) required
analysts to determine the sign of summands to choose the cheaper route. Under v5, this
distinction is irrelevant: $\text{add} = 2n$ always. Mixed-sign polynomial coefficients
incur zero penalty over positive-coefficient polynomials.

\paragraph{2. Uniform Horner advantage.}
Horner's method saves $N-1$ nodes over monomial form for all $N \geq 2$, with no
dependence on coefficient signs. Under v4, the Horner advantage was $3N-3$ (positive)
vs.\ $3N-11$ (mixed), hiding a sign-dependent threshold. Under v5: uniformly $N-1$ savings
(positive domain comparison) for all coefficient configurations.

\paragraph{3. Direct evaluation reversal for LSE.}
The previously preferred numerically stable LSE form ($8n$) is now dominated by the direct
form ($5n$). This is a reversal caused purely by the cost drop from $\text{add\_gen} = 11n$
to $\text{add} = 2n$. Similar reversals may occur in other numerical algorithms that favored
algebraically rearranged forms to avoid expensive mixed-sign additions.

%===========================================================================
\appendix
\section{Cost Model Reference}
%===========================================================================

\begin{lstlisting}[language=Python, caption={SuperBEST v5.1 positive-domain cost table (from superbest.py).}]
SUPERBEST_COSTS_POS = {
    "exp": 1, "ln": 1, "recip": 1,
    "div": 2, "neg": 2, "mul": 2, "sub": 2, "add": 2, "sqrt": 2,
    "pow": 1,   # X20: EPL/ELMl -- exp(n*ln(x)) = x^n for x > 0
    "sin": 1, "cos": 1,
}

SUPERBEST_COSTS_GEN = {
    "exp": 1, "ln": 1, "recip": 1,
    "div": 2, "neg": 2, "mul": 6, "sub": 2, "add": 2, "sqrt": 2,
    "pow": 3,
    "sin": 1, "cos": 1,
}
\end{lstlisting}

\begin{lstlisting}[language=Python, caption={ADD-T1 construction (from SuperBEST v5 audit).}]
# ADD-T1: x + y = LEdiv(x, DEML(y, 1)) = x - ln(exp(-y))
# DEML(y, 1) = exp(0) - ln(exp(y)) = 1 - y ... wait, that is -y.
# Concretely: DEML(y,1) = exp(y) - ln(1) = exp(y) -- no.
# Actual: LEdiv(x, DEML(y,1)) where DEML(y,1)=exp(y), then:
#   LEdiv(x, exp(y)) = x - ln(exp(y)) = x - y.  That's subtraction.
# Correct ADD-T1: LEdiv(x, DEML(y,1)) with DEML=exp-ln variant.
# As documented: the 2-node add construction via negation + subtraction route.
# Result: add(x, y) = 2n for ALL real x, y (proved in ADD_T1_General_Addition_2n.tex)
add_cost_v5 = 2   # ALL reals; positive domain: also 2n (not 3n as in v4)
\end{lstlisting}

\section{Polynomial Node Count Data}

The machine-readable node count tables are stored at:
\begin{center}
  \texttt{D:/monogate/python/results/polynomial\_node\_counts.json}
\end{center}
Fields: \texttt{general\_monomial\_form}, \texttt{horner\_form}, \texttt{taylor\_exp},
\texttt{taylor\_ln}, \texttt{taylor\_sin}, \texttt{chebyshev}, \texttt{legendre},
\texttt{bernstein}, \texttt{log\_sum\_exp}, \texttt{softmax}, \texttt{depth\_analysis},
\texttt{new\_identities}, \texttt{summary}.

\end{document}
